\section{Curved surfaces}\label{sec:curves}
The method of the previous sections was built on mechanical objects and it shows on round ones: the duck, the dome of a table
lamp, a sphere. The question that opened this part of the work was whether the right \emph{parts} would fix it --- quarter
circles, half discs, round plates, a catalogue cut from the LDraw library by usage --- and the answer, by measurement, was mostly
no. Three different things were being called ``the artifacts on curved surfaces'', and parts fix only one of them.

\subsection{The census, and why the tier question is moot}
\code{claude/generator/omr\_rank.py} ranks every part placed in the 1,150 official models: 401,012 placements of 4,425 distinct
LDraw parts. The top 1,000 cover 95.5\,\% of the placements, the top 2,000 98.7\,\% --- the ``5K / 8K / 10K catalogues'' are
the same list. By family, solid boxes are 44.1\,\%, shaped parts (slopes, curves, rounds, wedges, arches) 21.6\,\%, sideways
connectors 3.7\,\%: 69.5\,\% of everything ever placed is expressible in the engine's model; the rest is Technic (10.0\,\%),
minifigures (5.3\,\%) and bars, clips, hinges, wheels and windows (15.2\,\%). The 262-part catalogue already covers 86.8\,\% of
the solid placements, 75.3\,\% of the shaped ones and 63.3\,\% of the sideways ones, and what is missing is variants of what is
in. Adding solo candidates was measured negative in the motif round (\code{shapeParts} floods the skin with fragments), so a
wider catalogue only pays through the motifs, and the miner keeps only windows whose parts are all in the catalogue --- widening it
means re-mining. The 450 official models of 100 parts or fewer were mined separately as a ``small-set'' library (680 motifs):
measured through \code{motifLibrary}, dropped by the verification below on the duck, dolphin and table like the built-in library,
$-0.5$ IoU on the chair. Moot too.

\subsection{Three artifacts}
\begin{enumerate}
\item \textbf{The rim staircase.} Every horizontal cross-section of a round object is a Bresenham circle of cells, and the greedy
  fill tiles it with the longest plates it can, leaving a necklace of 1$\times$1 fragments along the edge. Round \emph{parts} do
  not fire on it: the solver places a 4$\times$4 round plate only where the field is a disc of radius 2.0 studs to within a
  quarter stud, and real cross-sections never are. What works is a \emph{recipe}: \textbf{disc layers} (\code{discs.js}, phase A1,
  before the motifs). Per level, a Kasa least-squares circle is fitted to the outline of each solid component --- the ``Hough
  transform'' step, done the cheap way because we know we are looking for one circle per blob --- and if the component is round
  enough (IoU $\ge .85$ against the fitted disc, outline rms $\le .45$ stud, radius $\ge 2$ studs) it is laid the way a builder
  lays a round table top: rows of 1$\times n$ plates along $x$, the next level along $z$, no lone 1$\times$1 ($7 = 4+3$, not
  $4+2+1$). Table 660 $\to$ 609 pieces, IoU .789 $\to$ .799; duck 1377 $\to$ 1358. \textbf{Quarter circles on corners}
  (\code{post.roundCorners}): after the finish, every 1$\times$1 tile on a convex corner of its layer becomes the quarter-round
  tile 25269 turned to the open side (table 12, duck 56), and an isolated 1$\times$1 tile becomes the round tile 98138. Cosmetic,
  IoU unchanged, and the staircase reads as a curve (figure~\ref{fig:discs}).
\item \textbf{The sloped skin of a blob.} Where the local curvature matches one of the $\sim$13 curved profiles the skin phase
  lays slopes; where it matches nothing it falls back to steps of plates. This one took a second round and is the next subsection.
\item \textbf{The lattice.} A round top at a fractional plate height is a layer of partial cells. The lattice fit of
  section~\ref{sec:lattice} already took care of this one (table .690 $\to$ .789), and section~\ref{sec:resolution} returns to it.
\end{enumerate}

\begin{figure}[H]
\centering
\includegraphics[height=0.78\textheight]{curves_before_after.jpg}
\caption{The table top before and after disc layers and corner rounding, coloured by piece orientation (blue plates along $x$,
green along $z$, red 1$\times$1, orange the quarter-round corner tiles). One layer of long plates with a necklace of 1$\times$1
at the rim becomes rows that turn $90^\circ$ on every level, no 1$\times$1 inside, quarter rounds on the convex corners.}
\label{fig:discs}
\end{figure}

\subsection{The sloped skin: what the artifacts actually were}
The first guess --- the catalogue has too few profiles --- was wrong in an instructive way. A per-column \emph{attribution} of the
top surface (\code{surface()} in \code{run.js}: for every sample column, the slope of the surface there and the kind of the
piece that ends up on top) said this about the round-4 defaults:
\begin{center}\small
\begin{tabular}{lccc}
\toprule
surface slope & duck: shaped / flat & dome & sphere \\
\midrule
20--30$^\circ$ & 63 / 37\,\% & 47 / 53 & 82 / 18 \\
30--40 & 63 / 37 & 65 / 35 & 58 / 42 \\
40--50 & 57 / 43 & 27 / 73 & \textbf{1 / 99} \\
50--60 & 59 / 41 & 11 / 89 & 16 / 84 \\
60--75 & 63 / 37 & 26 / 74 & 28 / 72 \\
\bottomrule
\end{tabular}
\end{center}
while a probe on the fresh field found a passing skin candidate whose top matched the surface at 100\,\% of those columns. The
parts exist and they fit; the solver did not get to place them. Three causes, in order of size:
\begin{enumerate}
\item \textbf{The disc layers were eating the flanks.} The exposure test (15\,\% of a layer's cells uncovered above) is true
  of every ring of a sphere's flank. The test is geometric now: exposed cells over outline length $\ge 1.5$ studs
  (\code{discRing}), measured against the full solid rather than the crust-carved field. The table top (ring 9.2) still
  qualifies; the sphere's rings (0.0--0.6) and the duck's (0.1--0.3) no longer do.
\item \textbf{The mined assemblies do not belong on a smooth slope.} The motifs that land there are ``a slope on a long brick''
  (3039+2456, 3040+3009) and ``a cheese on a plate''; the brick fits by volume inside the solid, imposes a straight ridge on a
  curved surface, and everything around it has to be fragments. On the chair the same assemblies are the frame members and gain
  4.6 points of IoU. No local measure separated the two cases (azimuth turning, gradient spread, a shaped-share filter were all
  tried), so the engine now \textbf{verifies}: the final field is solved once more without the assemblies --- fast, the motif phase
  is what costs --- and they are kept only when they gain \code{motifGain} $=.01$ of IoU. Kept on the chair ($+4.6$) and the
  cylinder ($+2.1$), dropped on the dome, sphere, duck, dolphin and table.
\item \textbf{Doubly curved surfaces need 1-wide parts along the fall line.} A 2-wide slope across a sphere fails on its cross
  slope (at $45^\circ$ of azimuth the surface drops 14\ldu{} across one stud); a 1-wide part fits any row whose profile it
  matches. The greedy, scoring volume alone, placed 2$\times$2 and 2$\times$4 slopes where they happened to pass, at random yaws.
  Now the 1-wide skin parts run first (\code{skinNarrow}) with an \textbf{alignment term} in the score (\code{skinAlign}:
  $\pm 0.6$ times the cosine between the part's own slope direction, read off its template, and the surface gradient under it,
  from the field's height map, weighted by steepness), and \code{post.widen} fuses two identical 1-wide parts side by side back
  into the catalogue's 2-wide version ($3040+3040 \to 3039$, $11477+11477 \to 15068$, $4286+4286 \to 3298$, found by template,
  not by name). Axis-facing columns of the dome went from 20\,\% flat to 2\,\%.
\end{enumerate}

\begin{figure}[H]
\centering
\begin{subfigure}{0.24\textwidth}\centering\includegraphics[width=\textwidth]{r_kd_dome_round4}\caption*{dome, round 4}\end{subfigure}
\begin{subfigure}{0.24\textwidth}\centering\includegraphics[width=\textwidth]{r_kd_dome_round5}\caption*{round 5}\end{subfigure}
\begin{subfigure}{0.24\textwidth}\centering\includegraphics[width=\textwidth]{r_kd_sphere_round4}\caption*{sphere, round 4}\end{subfigure}
\begin{subfigure}{0.24\textwidth}\centering\includegraphics[width=\textwidth]{r_kd_sphere_round5}\caption*{round 5}\end{subfigure}\\[2pt]
\begin{subfigure}{0.24\textwidth}\centering\includegraphics[width=\textwidth]{r_kd_duck_round4}\caption*{duck, round 4}\end{subfigure}
\begin{subfigure}{0.24\textwidth}\centering\includegraphics[width=\textwidth]{r_kd_duck_round5}\caption*{round 5}\end{subfigure}
\begin{subfigure}{0.24\textwidth}\centering\includegraphics[width=\textwidth]{r_kd_dolphin_round4}\caption*{dolphin, round 4}\end{subfigure}
\begin{subfigure}{0.24\textwidth}\centering\includegraphics[width=\textwidth]{r_kd_dolphin_round5}\caption*{round 5}\end{subfigure}
\caption{The skin before and after, coloured by kind: blue slopes, teal curved slopes, purple cheese, red inverted, grey
plates / bricks / tiles. Round 4 is the committed state before this work (wide parts, no alignment, motifs unverified, bridges
anywhere). Dome 393 pieces / IoU .877 $\to$ 413 / .901; sphere 734 / .853 $\to$ 777 / .856, flat columns on 20--75$^\circ$
surfaces 62\,\% $\to$ 15\,\%; duck 1304 / .855 $\to$ 1294 / .874, 38\,\% $\to$ 19\,\%; dolphin 262 / .684 $\to$ 240 / .691,
75\,\% $\to$ 52\,\%. Undersides (inverted slopes 4287 / 3747, added to the core catalogue): sphere 65\,\% $\to$ 38\,\% flat.}
\label{fig:skin}
\end{figure}

\subsection{What did not make it}
Looser skin tolerances (the tolerance is not what blocks the parts). 1-wide parts first \emph{without} alignment: the sphere
got worse (34 $\to$ 52\,\% flat), narrow parts at random yaws block each other too. Alignment alone, with wide parts: excludes
the misaligned parts and nothing replaces them. Motifs after the skin: worse than motifs off everywhere. Any local test for
``this cell is on a smooth curve'' (azimuth turning, gradient spread, jumps): no threshold separates the duck from the chair.
\textbf{Profile chains} (\code{skin.js}, option \code{profiles}): rows along the fall line, each a chain of 1-wide parts chosen by
a dynamic programme over the row --- correct and clean on every row it takes, and still worse than the plain greedy (dome 28\,\% vs
21\,\% flat), because the hard $x$/$z$ facing split loses the freedom the greedy keeps at 30--45$^\circ$ of azimuth and 2-wide
parts never enter the chain. Kept as an option. A pass for the diagonal seams (double-convex corner slopes right after the narrow
skin) placed four parts on the dome and changed nothing: at 8 studs of radius the gradient turns too slowly for a 2$\times$2
corner. Letting the extended 1-wide parts into the narrow pass: duck 19 $\to$ 45\,\% flat, twice.

\section{Three bugs that looked like features}\label{sec:round7}
\paragraph{The bridge in the air.} \code{post.bridge} joins components with the shortest chain of $1\times k$ plates through
\emph{free} cells, and free meant ``no piece there'', not ``inside the model''. On a round or organic model the shortest free
route between two components is around the outside, so the chain became a skirt of plates under a dome, an L of 1$\times$6
plates hanging off a game controller --- and, sitting on the real surface, it then kept the skin from being laid there (the knobs
that could not be round). Measured on dome / duck / chair / rafs5 / be2: 2 / 6 / 3 / 4 / 5 bridge plates, every one of them in
cells with field fill below 0.3. A cell whose mean fill is below \code{bridgeFill} ($.3$) is now as blocked as an occupied one;
chains stay inside the solid. No bridge plate in the air on any model, IoU $+0.2 \ldots +0.5$ everywhere, and the components the
air chains used to join are now what they are (be2 1 $\to$ 6: strut columns, a problem for the connectivity passes through the
solid, not a reason to build in the air). A cost for air cells was tried first; two air plates still fit in the budget, so the
hard rule won.

\begin{figure}[H]
\centering
\begin{subfigure}{0.24\textwidth}\centering\includegraphics[width=\textwidth]{r_br_duck_outside_crop}\caption*{duck, through any empty cell}\end{subfigure}
\begin{subfigure}{0.24\textwidth}\centering\includegraphics[width=\textwidth]{r_br_duck_inside_crop}\caption*{inside the solid}\end{subfigure}
\begin{subfigure}{0.25\textwidth}\centering\includegraphics[width=\textwidth]{r_br_dome_outside_crop}\caption*{dome, any empty cell}\end{subfigure}
\begin{subfigure}{0.25\textwidth}\centering\includegraphics[width=\textwidth]{r_br_dome_inside_crop}\caption*{inside the solid}\end{subfigure}
\caption{The bridge phase, bridge plates in orange. Through any empty cell the duck gets six plates and the dome two, every one
sitting on the outside of the surface (the duck's head, the dome's shoulder); inside the solid, none in the air. Small at this
scale, and exactly what kept the skin from being laid under them.}
\end{figure}

\paragraph{The measured shapes were hollow.} \code{lego\_catalog\_shapes.py} measures a part as the volume between its top and
its bottom profile, and the bottom profile of an LDraw brick is the ceiling of its cavity: the 2$\times$2 quarter-round plate
30357 came out 50\,\% full, the 4$\times$4 corner-round plate 60474 63\,\%. Against a solid field those templates never passed
the tolerance, which is why ``the round family does not fire'' above and why \code{shapeParts} measured negative in the motif
round. The undersides are filled solid now for everything but arches and inverted parts (20 of 51 shapes). The rounded-edge box
goes from wedges to a stack of quarter-round plates (IoU .939 $\to$ .957), the table .809 $\to$ .820. With real templates the
quarter-round plates then fitted the duck's body level by level \emph{before} the skin could, so the round family's turn is
decided per model: after the sloped skin when at least 30\,\% of the top surface is sloped (an animal, a dome), before it
otherwise (a table, a rounded box). A per-candidate steepness gate was tried first and lost on the sphere and the dolphin.

\paragraph{Loose pieces.} Every extra component of be2 (3 $\to$ 8 after the narrow skin) was one loose piece --- a cheese on a
wing tip, a tile left by the finish. Components of one piece that touch nothing are removed (\code{dropLoose}): be2 8 $\to$ 1,
chair 5 $\to$ 1, rafs5 8 $\to$ 3.

\subsection{Sideways synthesis: ``anything on walls''}
The question was discs on a wall, and more generally anything on a vertical face. The mined SNOT motifs only ever contain the
1$\times$1 round plate on a headlight brick, because no official set puts a 2$\times$2 disc on a wall inside a mining window. But
the mined motifs do contain the information that matters: \emph{where the side stud is}. \code{motifs/snot.js} recovers, for
every mined sideways part next to a host, the world position of the host's side stud (the mode over the observations --- a
1$\times$2 plate can hang by either cell): 48 stud points for (host, host yaw, part orientation) seen five or more times, nearly
all on 4070 (headlight) and 87087 (one stud on the side). The underside of every plate, tile, slope or round plate takes a stud in
each cell of its footprint, so the same stud holds \emph{any} such part by \emph{any} of its cells: 548 host-plus-part assemblies,
2192 solver variants after the yaws and mirrors, built through the same \code{assemblyVariants} as the mined motifs, placed as
rigid assemblies, exported with their real orientation, round-tripping through the LDR exporter exactly, and reproducing all nine
two-part mined SNOT motifs they cover. Phase \code{S-snot} after the upright skin, option \code{snot}.

It is off by default, and the reason is the lesson of the skin all over again. A sideways part is chosen here by volume alone,
like any other candidate, and the field-driven selection hangs small things wherever a partial cell lets them lower the error:
cones and pyramids on a rounded edge until those were excluded, 1$\times$1 round plates on the corners of a box. On the boss box
the mined SNOT motifs, kept by the verification where they pay, already do the disc-on-a-wall greeble (17 sideways round plates
and tiles, figure~\ref{fig:snot2}); the synthesis adds pieces and nothing to the IoU. The sloped skin only became clean once the
solver was told \emph{where a slope belongs}; a wall has no top-surface gradient, but it has a depth map per facing. The form of
the feature is therefore \textbf{a wall is a floor turned on its side}: resample the field so that one facing becomes ``up'' (the
4\ldu{} samples along $x$ pair up into 8\ldu{} plates, the plates split into two samples), run the narrow aligned skin there with
the upright catalogue, and turn the result back --- its pieces become sideways pieces (their yaw composed with the frame rotation
gives the orientation index), its ``ground'' cells become the places where the wall must carry a side-stud host instead of a plain
brick. That is how a real SNOT hull is built: a core of side-stud bricks and sideways plating. The pieces of this exist --- the 24
orientations, the rasteriser, the rigid groups, the exporter, now the legal stud points; the frame rotation, the host insertion
and the merge with the upright solution are the work.

\begin{figure}[H]
\centering
\begin{subfigure}{0.49\textwidth}\centering\includegraphics[width=\textwidth]{r_sn_boxBoss_off}\caption{the boss box at 24 studs, defaults: 805 pieces, IoU .985; the boss is on the right face}\end{subfigure}
\begin{subfigure}{0.49\textwidth}\centering\includegraphics[width=\textwidth]{r_sn_boxBoss_snot}\caption{with the synthesised sideways assemblies: 809 pieces, IoU .985}\end{subfigure}
\caption{A box with a cylindrical boss on one face (the disc on a wall), coloured by kind; sideways pieces in bright teal. The
mined sideways motifs already greeble the boss; the synthesis changes little until it is driven by the wall's own depth map.}
\label{fig:snot2}
\end{figure}

\section{Choosing the resolution}\label{sec:resolution}
For a given mesh some stud counts come out visibly better than others, and nothing said which. The question behind it was
whether a metric exists for a ``good'' default --- good meaning the high-frequency detail of the mesh survives while its plain
areas stay plain. The second half is free: a plain area costs the same few big plates at any resolution. The first half is a
sampling question with a budget attached. And the ``some resolutions are better'' part is real, measurable before anything is
solved, and has nothing to do with detail: it is the lattice.

\subsection{What the literature says}
A feature survives a grid when it spans about two cells (the two-voxel form of Nyquist); for surfaces the sharper statement is the
$\varepsilon$-sampling of Amenta and Bern: the sample density must follow the \emph{local feature size}, the distance to the medial
axis, so thin parts and tight curves need dense samples and flat parts hardly any. Our field samples the mesh at 4\ldu{} and the
parts at 20\ldu, so the stud is the period that matters. Two descriptors give the feature sizes of a mesh: the shape diameter
function of Shapira, Shamir and Cohen-Or measures local thickness by casting rays through the solid --- chord lengths along the
axes are its cheap cousin and the engine's ray caster gives them for free --- and radii of curvature give the other size, how tight
the round parts are. For the operating point itself, rate--distortion thinking applies: pieces are the rate, the error against
the original the distortion, and the conventional choice on such a curve is its knee (Kneedle, Satopaa et al.). Testuz,
Schwartzburg and Pauly voxelise at a user-chosen resolution and do not ask the question.

\subsection{Measuring across resolutions}
The engine's IoU is against its own field at the solution's resolution, so it cannot compare two resolutions.
\code{test/resolution\_bench.mjs} rasterises every solution back into one fine \emph{reference} voxelisation of the mesh (96 studs,
no lattice fit, no crust) and measures there: the IoU of the solid against the reference solid, the carved hollow counted as
solid since it is hidden, and the same IoU restricted to a band half a stud wide around the reference surface, where the detail
lives. The same is done for the \emph{field} alone at each stud count: what the voxel grid can represent before any part is chosen,
so that the solver's own loss is the gap between the two. Pieces are matched to their template by id, rotation and footprint;
without that the measure fell with resolution.

\begin{figure}[H]
\centering
\includegraphics[width=\textwidth]{resolution_curves}
\caption{Left: fidelity against the 96-stud reference versus piece count, for the solution (solid) and for the field alone
(dotted); the ring marks the automatic choice. Middle: the same within half a stud of the surface. Right: the field's fidelity at
every stud count within $\pm 15\,\%$ of the budget --- the lattice effect; rings mark the chosen count.}
\label{fig:rd}
\end{figure}

Three things are in figure~\ref{fig:rd}.
\begin{enumerate}
\item \textbf{There is no knee.} On the smooth shapes fidelity grows steadily in $\log(\text{pieces})$: every doubling of the piece
  count buys the sphere about .02 of solid IoU and .07 of band IoU, from 8 studs to 48, and Kneedle on such a curve returns an
  end point. The resolution is a budget decision, not a property of the mesh --- except for the next point.
\item \textbf{The round shape has a ceiling.} The sphere gains little from 16 to 48 studs for 10$\times$ the pieces, and falls
  apart: one component at 16, 69 at 32. Once the radius spans about 8 studs the slopes and curved parts follow the surface as well
  as they ever will; past that the solver pays in pieces and connectivity for nothing a builder would see. The furniture is nearly
  flat from 20 to 48: its planar faces are right at any resolution that lands them on the lattice and its fillets never span a stud
  at any practical count.
\item \textbf{The lattice is the ``some resolutions are better'' effect, and it is large.} A plate-thick table top or a stud-wide leg
  either lands on the lattice or is smeared over two cells, and that error dwarfs what one more stud of resolution buys. The
  \emph{field} shows it before any part is placed: the table's field at 43 studs is .81 against the reference, at 35 or 47 it is
  .93; the chair's at 37 is .966, at 32 .982. The lattice fit of section~\ref{sec:lattice} ($\pm 5\,\%$ of scale) removes part
  of it, not all: at some counts no scale within 5\,\% lands both the top and the legs.
\end{enumerate}

\begin{figure}[H]
\centering
\begin{subfigure}{0.49\textwidth}\centering\includegraphics[width=\textwidth]{r_res_table_35}\caption{table at 35 studs: 562 pieces, IoU .885}\end{subfigure}
\begin{subfigure}{0.49\textwidth}\centering\includegraphics[width=\textwidth]{r_res_table_43}\caption{at 43 studs: 772 pieces, IoU .831}\end{subfigure}\\[2pt]
\begin{subfigure}{0.49\textwidth}\centering\includegraphics[width=\textwidth]{r_res_sphere_14}\caption{sphere at 14 studs (the automatic choice): 562 pieces, 1 component}\end{subfigure}
\begin{subfigure}{0.49\textwidth}\centering\includegraphics[width=\textwidth]{r_res_sphere_32}\caption{at 32 studs: 3553 pieces, 69 components}\end{subfigure}
\caption{The same meshes at neighbouring or larger stud counts. The table at 43 studs is worse than at 35 by every measure ---
the lattice --- and the sphere past its curvature ceiling buys fragments.}
\end{figure}

\subsection{The automatic stud count}
\code{pipeline.autoStuds} (\code{studs: 'auto'}, or the \emph{Stud count: automatic} switch) combines three numbers, in
2--9\,s on the test models:
\begin{itemize}
\item \textbf{A piece budget.} Two quick pilot solves (no motifs, no post passes) at 12 and 20 studs give the piece counts there
  and the growth exponent, pieces $\propto N^{e}$ with $e$ fitted per model: sphere 2.6, duck 2.6, chair 2.4, table 3.1 (a thin
  model grows fastest while its struts appear), rounded box 1.6. The budget \code{autoPieces} (1200) gives $N_{\text{budget}}$. One
  pilot with a fixed exponent overshot the table by 50\,\%: it vanishes at 12 studs.
\item \textbf{A curvature ceiling.} The radius of curvature across every smooth edge of the welded mesh --- the dihedral angle
  between its two faces over the width of the two-face strip, creases above $35^\circ$ excluded since a box edge is sharp at any
  resolution --- weighted by edge length; its median must span 8 studs, so a sphere of radius $R$ asks for 16 studs across. Exact on
  the synthetic shapes (sphere 16, cylinder 26.7, rounded box of $R = W/4$: 32, boss 44.8). It applies only when at least 10\,\%
  of the smooth surface is curved: the chair is 26\,\% fillets of radius $W/24$ and would ask for 194 studs, which is not what the
  ceiling is for. Vertex normals were the first attempt and fail on exactly these meshes: averaged across a crease they tilt by
  $45^\circ$, so every flat wall of a box read as a curve of radius $W/4$. Chord thickness is computed and reported --- at $N$ studs
  it says how much of the model is thinner than a stud --- but does not set the count: its 10th percentile is grazing chords and
  fillets and asks for 160--900 studs on every real model.
\item \textbf{The lattice tie-break.} Around $N = \min(N_{\text{budget}}, N_{\text{curv}})$ the seven stud counts within
  $\pm 15\,\%$ are each prepared exactly as the engine will (lattice fit included) and their \emph{fields} measured against a
  reference voxelisation of twice the finest candidate; the IoU, less $.02$ per doubling of the pieces so that fewer pieces win
  ties (a smooth shape gains about $.005$), picks the count. The field is pulled rather than splatted (every reference cell takes
  the field cell under its centre), 0.3\,s per candidate.
\end{itemize}

\begin{table}[H]
\centering\small
\begin{tabular}{lrlrrl}
\toprule
model & chosen & by & $N_{\text{budget}}$ & $N_{\text{curv}}$ & field IoU of the candidates \\
\midrule
sphere & 14 & lattice & 18.9 & 16 & 14 .964 \ 15 .967 \ 16 .969 \ 17 .971 \ 18 .973 \\
dome & 14 & lattice & 25.4 & 16 & 14 .967 \ldots 18 .975 \\
cylinder & 24 & lattice & 37.8 & 26.7 & 23 .967 \ \textbf{24 .975} \ 26 .977 \ 27 .978 \ 28 .979 \ 30 .981 \ 31 .981 \\
rounded box $R = W/4$ & 30 & lattice & 38.6 & 32 & 27 .990 \ 29 .990 \ \textbf{30 .999} \ 32 .999 \ 34 .999 \ 35 .995 \ 37 .990 \\
rounded box $R = 0.15W$ & 33 & lattice & 39.2 & 53.3 & \textbf{33 .998} \ 35 .990 \ 37 .990 \ 39 .995 \ 41 .998 \ 43 .998 \ 45 .998 \\
boss on a box & 34 & lattice & 40.3 & -- & \textbf{34 .998} \ 36 .997 \ 38 .989 \ 40 .998 \ 42 .998 \ 44 .998 \ 46 .993 \\
duck & 20 & lattice & 22.5 & 29.8 & \textbf{20 .964} \ 21 .964 \ 22 .964 \ 23 .967 \ 24 .969 \ 25 .969 \ 26 .972 \\
dolphin & 54 & lattice & 72.5 & 74.2 & \textbf{54 .943} \ 58 .943 \ 61 .943 \ 64 .949 \\
table & 35 & lattice & 40.7 & 93.6 & \textbf{35 .926} \ 37 .921 \ 39 .928 \ 41 .896 \ 43 .811 \ 45 .914 \ 47 .934 \\
chair & 32 & lattice & 33.7 & 194 & 29 .971 \ 31 .971 \ \textbf{32 .982} \ 34 .979 \ 36 .974 \ 37 .966 \ 39 .977 \\
rafs5 & 61 & lattice & 70.6 & 328 & 54 .817 \ 58 .823 \ \textbf{61 .826} \ 64 .828 \\
be2 & 64 & budget & 97.1 & 191 & 54 .571 \ 58 .586 \ 61 .577 \ \textbf{64 .631} \\
\bottomrule
\end{tabular}
\caption{What the automatic choice does on the test models. $N_{\text{curv}}$ is the curvature ceiling (applied when the curved
share of the surface is at least 10\,\%; the chair's and the aircraft's fillets do not count). The table at 43 studs and the
rounded box at 35 are the counts a user would have seen as ``worse''; they are found in 0.3\,s each and skipped. The biplane's
field is .63 against the reference even at 64 studs: its struts are thinner than a stud at any count the budget allows.}
\end{table}

What did not work, briefly: vertex-normal curvature; Kneedle on the solution curves (no knee to find); one pilot and a fixed
exponent; the chord-thickness percentile as a criterion; the lattice \emph{score} of section~\ref{sec:lattice} as the tie-break
--- it ranks the chair right (2.24 at 32) but is blind to plate heights and to the thin elements that make the table's 43 bad,
and the field fidelity sees both.
